\subsection{A Future Is a Reply Stream}
\label{sec:futures}

The supplied future library can collect more than one reply from a request.
The send syntax uses a dollar sign:
\code{[receiver <= message \$ future-object]}.
Each \code{!} in the receiver's handler appends another reply to the future.
This is more general than a promise that resolves to one value.

To run the following example from the repository root:

\begin{lstlisting}
(inc "tutorial/examples/04-futures" :use-source t)
[range-client <== [:collect 1 4]]
; => (1 2 3 4)
\end{lstlisting}

The file loads the supplied future library. Here is the server:

\wholelisting[firstline=4,lastline=11]{tutorial/examples/04-futures.abclr2}

The server sends each integer and then \code{:done}.
The terminal message is part of our protocol: futures themselves do not
know that a producer has finished.

\wholelisting[firstline=12]{tutorial/examples/04-futures.abclr2}

The initial send is asynchronous. Each \code{next-value} waits if no reply
is currently available. Thus waiting occurs at the point where the client
needs a particular element, rather than at the initial request.
The sentinel is removed from the stream but excluded from the result.
The range from 3 to 2 is empty and returns \code{nil}; the server still
sends \code{:done}, so the client does not wait forever.

\subsection{Peeking, Consuming, and Readiness}

\begin{longtable}{L{0.39\textwidth}L{0.51\textwidth}}
\textbf{Operation} & \textbf{Behavior in the supplied implementation}\\\hline
\endhead
\code{(make-future)} & Create and initialize an empty future.\\
\code{(ready? f)} & Return whether at least one reply is buffered.\\
\code{(next-value f)} & Wait for and remove the next reply.\\
\code{(next-value f :remove nil)} & Wait for the next reply without removing it.\\
\code{(all-values f :wait nil)} & Return the current buffered replies,
or \code{nil} if none exist; removal defaults to true.\\
\code{(reset-future f)} & Send a reset request; use only when old producers
and readers have finished.\\
\end{longtable}

The actual macros in \filepath{src/user.lsp} default \code{:remove}
to true for \code{next-value} and \code{all-values}; an old comment
in the library header gives a different default.
\code{all-values} is a snapshot of buffered values, not an operation
that waits for a terminal marker. Its \code{:wait t} default waits
when the buffer is empty, but does not establish producer completion.

For stream consumers, use \code{next-value} and an explicit terminal
marker. Use \code{all-values} with \code{:wait nil} to inspect the replies
currently buffered; a snapshot is not a producer-completion signal.

\begin{codeblock}
\begin{lstlisting}
(let ((f (make-future)))
  [range-server <= [:range 8 8] $ f]
  (list (next-value f :remove nil)
        (next-value f)
        (next-value f)
        (ready? f)))
; => (8 8 :DONE NIL)
\end{lstlisting}
\end{codeblock}

The implementation records one waiting consumer destination.
Treat each future as a stream with one active consumer; use a separate
future per request rather than sharing one among unrelated readers.

\paragraph{Exercise 4.}
Change the server to send squares, keeping the terminal marker.
Then change the payload to \code{[:value n]} so that a legitimate data
value cannot be confused with \code{:done}.
Compare the two message shapes with Section~\ref{sol:stream}.
